\newpage
\section{Introduction}
\label{ch:introduction}

\subsection{Cloud Computing and the Rise of Container Orchestration}

The rapid evolution of modern software systems has fundamentally transformed the way
infrastructure is designed, deployed, and operated. At the heart of this transformation
lies \textit{cloud computing} — a paradigm that enables on-demand access to a shared pool
of configurable computing resources, including networks, servers, storage, applications,
and services, that can be rapidly provisioned and released with minimal management effort
(\cite{nist:cloudcomputing}). This model has shifted the industry away from static,
on-premises data centres toward elastic, scalable environments that can react dynamically
to workload demands.

As organisations increasingly adopt cloud-native architectures, the notion of designing
systems that are resilient, scalable, and operationally excellent has become central to
engineering practice. The AWS Well-Architected Framework, for instance, codifies these
concerns into a set of established design principles for container-based workloads,
emphasising repeatability, automation, and observability as first-class concerns in
modern infrastructure (\cite{aws:containerlens}).

A critical enabler of cloud-native architectures is the \textit{container} — a
lightweight, self-contained unit of software that packages application code together with
its dependencies, configuration, and runtime environment. Unlike traditional virtual
machines, containers share the host operating system kernel and therefore incur
significantly less overhead while still providing process-level isolation
(\cite{docker:containersvsvms}). This characteristic makes containers especially well
suited to microservice architectures, where large monolithic applications are decomposed
into smaller, independently deployable services that can be scaled individually based on
demand.

\subsection{Kubernetes as the Standard for Container Orchestration}

Managing containerised workloads at scale introduces significant operational complexity.
Scheduling containers across a fleet of machines, maintaining desired application state,
performing rolling updates, and recovering from failures all require robust automation.
It is in response to exactly these challenges that \textit{Kubernetes} — also referred
to as K8s — emerged as the dominant container orchestration platform.

Originally designed and used internally at Google as the \textit{Borg} system, Kubernetes
was open-sourced in 2014 and donated to the Cloud Native Computing Foundation (CNCF),
where it has since grown into the second-largest open-source project in the world
(\cite{ibm:kuberneteshistory}). Kubernetes provides a declarative API for defining the
desired state of an application cluster; its control plane continuously reconciles the
observed state of the system against this desired state, automatically scheduling pods,
managing service discovery, and orchestrating storage (\cite{kubernetes:overview}).

To reduce operational complexity further, lightweight Kubernetes distributions such as
\textit{K3s} — developed by Rancher and certified by the CNCF — have gained widespread
adoption. K3s packages the full Kubernetes feature set into a single binary with a
significantly reduced memory footprint, making it suitable for edge computing, IoT
environments, and resource-constrained infrastructure (\cite{k3s:docs}). In the context
of this thesis, K3s serves as the underlying cluster management layer for the
experimental infrastructure under investigation.

\subsection{The Challenge of Scaling Stateful Applications}

While Kubernetes excels at managing stateless workloads — where any replica of a
service can handle any request independently — scaling \textit{stateful applications}
presents a distinct set of challenges. Stateful applications, such as in-memory databases
and message queues, maintain persistent data or internal state that must be carefully
managed across the lifecycle of individual pods. Kubernetes provides the
\texttt{StatefulSet} primitive to address these requirements, offering stable network
identities, ordered deployment and scaling, and persistent volume binding
(\cite{kubernetes:statefulsets}).

\textit{Redis} is a canonical example of such a stateful application. It is a
high-performance, in-memory data store widely used for caching, session management, and
real-time analytics (\cite{redis:io}). When deployed as a Redis Cluster across multiple
nodes, it uses sharding to distribute data and replication to ensure fault tolerance.
Scaling a Redis Cluster in a Kubernetes environment is non-trivial: adding or removing
nodes requires coordinated resharding of data slots, careful management of cluster
topology, and maintenance of quorum — operations that introduce significant risk if
performed reactively under resource pressure.

Recent work from the CNCF community has highlighted these concerns, noting that while
Kubernetes operators can partially automate lifecycle management for Redis, the path
toward robust, cloud-native stateful services remains an open challenge
(\cite{cncf:redisoperator}).

\subsection{Autoscaling Strategies: From Reactive to Predictive}

Kubernetes provides native support for horizontal scaling through the
\textit{Horizontal Pod Autoscaler} (HPA), which monitors CPU and memory utilisation
metrics and adjusts the number of pod replicas accordingly (\cite{kubernetes:autoscaling}).
While the HPA is effective for stateless services under gradually increasing load, its
reactive nature introduces an inherent latency: scaling decisions are made only
\textit{after} resource thresholds have already been breached. For latency-sensitive or
stateful workloads, this lag can result in degraded performance, pod restarts, or
cascading failures during traffic spikes (\cite{thoras:hpa}).

An increasingly popular alternative is \textit{Kubernetes Event-Driven Autoscaling}
(KEDA), a CNCF-graduated project that extends the standard HPA to support event- and
metric-driven scaling from a rich ecosystem of external sources, including message
queues, databases, and custom Prometheus metrics (\cite{keda:home}; \cite{cncf:kedagrad}).
By decoupling scaling decisions from simple resource utilisation and enabling them to
respond to application-level events, KEDA enables a more proactive and context-aware
approach to scaling. This is particularly relevant for stateful workloads where
pre-emptive resource provisioning may prevent the cluster-level disruptions that reactive
scaling cannot avoid.

\subsection{Research Objective and Thesis Structure}

This thesis investigates the performance differences between two autoscaling strategies
for stateful applications deployed in a Kubernetes cluster managed by K3s. The first
strategy employs the native Horizontal Pod Autoscaler in a \textit{reactive} configuration,
scaling a Redis Cluster in response to CPU and memory resource expenditure. The second
strategy employs a \textit{custom KEDA-based operator} in a predictive, event-driven
configuration, triggering scaling decisions based on application-level metrics sourced
from a Prometheus observability stack.

The experimental infrastructure consists of a FastAPI-based entrypoint service
(\cite{fastapi:docs}), a Redis Cluster with a minimum of three nodes, an observability
stack comprising Prometheus and Grafana deployed via Helm, and two scaling
configuration manifests representing the naive HPA and KEDA-driven approaches
respectively. All components are deployed on a K3s cluster, providing a reproducible
and resource-constrained environment representative of real-world edge and on-premises
deployments.

The remainder of this thesis is structured as follows. Chapter ~\ref{sec:stateful-vs-stateless} details the differences between the automatic provisioning of stateless and stateful applications, a well as the issues that arise when scaling stateful workloads. Chapter~\ref{sec:technologies} provides the theoretical background on Kubernetes architecture, autoscaling mechanisms,
stateful workloads, and the Redis Cluster model. Chapter~\ref{sec:architecture} describes the
system design and experimental setup. Chapter~\ref{sec:results} presents the
experimental results and performance analysis. Chapter~\ref{sec:conclusion} summarises the
findings and discusses directions for future work.